\chapter{Materiales e Instrumental}
Para el desarrollo de las distintas experiencias de este trabajo práctico se emplearon los siguientes elementos e
instrumentos:

\begin{itemize}
    \item \textbf{Osciloscopio Digital:} UNI-T UTD2102CEX, ancho de banda de $100\text{ MHz}$, tasa de muestreo de
          $1\text{ GS/s}$. Utilizado con puntas de prueba atenuadoras $10\times$.
    \item \textbf{Multímetro Digital True RMS:} UNI-T UT61D+, utilizado para las mediciones de valor eficaz verdadero
          en señales sinusoidales y deformadas.
    \item \textbf{Multímetro Digital (Respuesta al Valor Medio):} UNI-T UT33A+, utilizado para contrastación y análisis
          del error sistemático por calibración.
    \item \textbf{Medidor Digital de Potencia y Factor de Potencia:} Vatímetro monofásico de tomacorriente (plug-in),
          utilizado para verificar el factor de potencia y la potencia activa del circuito dimmerizado.
    \item \textbf{Medidor LCR Digital:} Agilent U1733C, utilizado para la medición de la capacidad real de los
          capacitores de compensación.
    \item \textbf{Transformador de Aislación:} Relación 1:1 ($220\text{ V}$ a $220\text{ V}$), utilizado como elemento de
          protección para el osciloscopio y para independizar las masas de medición.
    \item \textbf{Módulo de Carga Reactiva:} Montaje experimental con tubo fluorescente común y balasto inductivo (reactancia
          magnética Philips).
    \item \textbf{Capacitores de Compensación:} Banco de capacitores de valores varios para la corrección del factor de
          potencia.
    \item \textbf{Módulo de Regulación de Fase:} Circuito dimmer con control por Triac y lámpara incandescente como carga.
\end{itemize}
